% ----- CHAUPAI 29 -----

\newpage

\subsection*{\deva{एकोनत्रिंशत्तम चौपाई} (Twenty-Ninth Chaupai)}
\addcontentsline{toc}{subsection}{\deva{एकोनत्रिंशत्तम चौपाई} (Twenty-Ninth Chaupai) — \deva{चारों जुग...}}

\begin{minipage}[t]{0.55\textwidth}
\vspace{0pt}

{\large \linespread{1.5}\selectfont
\deva{चारों जुग परताप तुम्हारा।}\\
\deva{है परसिद्ध जगत उजियारा॥}
\par}

\vspace{0.5em}

{\large \linespread{1.5}\selectfont
\telu{చారోఁ జుగ పరతాప తుమ్హారా।}\\
\telu{హై పరసిద్ధ జగత ఉజియారా॥}
\par}

\vspace{0.5em}

{\large \linespread{1.3}\selectfont
\textit{cārõ juga paratāpa tumhārā |}\\
\textit{hai parasiddha jagata ujiyārā ||}
\par}
\end{minipage}%
\hfill
\begin{minipage}[t]{0.42\textwidth}
\vspace{0pt}
\begin{tcolorbox}[anchorbox]
\textbf{Timeless Principles:} His glory shines in all four epochs (\textit{Charon Yug}). In the *Valmiki Ramayana* (*Sundara Kanda, 1:198*), the celestials praise Hanuman, declaring that anyone who possesses his four timeless qualities—courage, vision, intellect, and skill—will never fail in their endeavors. These core principles remain globally relevant across all generations.
\vspace{0.5em}\newline Hanuman's four foundational leadership qualities guarantee success that remains relevant across all space and time.\newline\null\hfill\href{https://www.valmiki.iitk.ac.in/sloka?field_kanda_tid=5&language=dv&field_sarga_value=1}{\textit{Sundara Kanda, 1:198}}
\end{tcolorbox}
\end{minipage}

\vspace{0.5em}
\subsubsection*{\textbf{Visual Rhythm Map:}}
\begin{table}[h!]
\centering
\renewcommand{\arraystretch}{1.2}
\setlength{\tabcolsep}{0pt}
\begin{tabularx}{\textwidth}{|*{16}{>{\centering\arraybackslash\hsize=1.0\hsize}X|}}
\hline
\multicolumn{4}{|>{\centering\arraybackslash\hsize=4.0\hsize}X|}{\textbf{\guru{चा}\guru{रों}} \par (4)} &
\multicolumn{2}{>{\centering\arraybackslash\hsize=2.0\hsize}X|}{\textbf{\laghu{जु}\laghu{ग}} \par (2)} &
\multicolumn{5}{>{\centering\arraybackslash\hsize=5.0\hsize}X|}{\textbf{\laghu{प}\laghu{र}\guru{ता}\laghu{प}} \par (5)} &
\multicolumn{5}{>{\centering\arraybackslash\hsize=5.0\hsize}X|}{\textbf{\laghu{तु}\guru{म्हा}\guru{रा}} \par (5)} \\
\hline
\end{tabularx}
\vspace{-1.5em}
\end{table}

\begin{table}[h!]
\centering
\renewcommand{\arraystretch}{1.2}
\setlength{\tabcolsep}{0pt}
\begin{tabularx}{\textwidth}{|*{16}{>{\centering\arraybackslash\hsize=1.0\hsize}X|}}
\hline
\multicolumn{2}{|>{\centering\arraybackslash\hsize=2.0\hsize}X|}{\textbf{\guru{है}} \par (2)} &
\multicolumn{5}{>{\centering\arraybackslash\hsize=5.0\hsize}X|}{\textbf{\laghu{प}\laghu{र}\guru{सि}\laghu{द्ध}} \par (5)} &
\multicolumn{3}{>{\centering\arraybackslash\hsize=3.0\hsize}X|}{\textbf{\laghu{ज}\laghu{ग}\laghu{त}} \par (3)} &
\multicolumn{6}{>{\centering\arraybackslash\hsize=6.0\hsize}X|}{\textbf{\laghu{उ}\laghu{जि}\guru{या}\guru{रा}} \par (6)} \\
\hline
\end{tabularx}
\vspace{-1.5em}
\end{table}

\vspace{1em}
\textbf{ \deva{सरल-संस्कृतम्}:}\\
 \deva{चतुर्षु युगेषु भवतः प्रतापः (महिमा) अस्ति।}\\
 \deva{संपूर्ण-जगति भवतः प्रकाशः प्रसिद्धः अस्ति॥}\\


Your glory and majesty exist throughout the four Yugas / ages.\\
Your illuminating light is famous across the universe.


\vspace{1em}
\newpage
\textbf{\deva{प्रतिपदार्थः} (Meaning of each word):}
\begin{table}[h!]
\centering
\renewcommand{\arraystretch}{1.2}
\begin{tabularx}{\textwidth}{@{} l l X X @{}}
\toprule
\textbf{Awadhi} & \textbf{Sanskrit} & \textbf{English} & \textbf{Grammar \& Notes} \\
\midrule
\deva{चारों जुग} \glossaryicon / \telu{చారోం జుగ} / \textit{cārõ juga}  &  \deva{चतुर्षु युगेषु}  &  In the four Yugas  & \\ 
\deva{परताप} \gotchaicon / \telu{పరతాప} / \textit{paratāpa}  &  \deva{प्रतापः}  &  Glory / Majesty  & Swara-Bhakti \\
\deva{तुम्हारा} / \telu{తుమ్హారా} / \textit{tumhārā}  &  \deva{भवतः}  &  Your  &   \\
\deva{है} \gotchaicon / \telu{హై} / \textit{hai}  &  \deva{अस्ति}  &  Is  &   \\
\deva{परसिद्ध} \gotchaicon / \telu{పరసిద్ధ} / \textit{parasiddha}  &  \deva{प्रसिद्धः}  &  Famous / Known  & Swara-Bhakti \\
\deva{जगत} / \telu{జగత} / \textit{jagata}  &  \deva{जगति}  &  In the world  &   \\
\deva{उजियारा} / \telu{ఉజియారా} / \textit{ujiyārā}  &  \deva{प्रकाशः / उज्ज्वलः}  &  Light / Radiance  &   \\
\bottomrule
\end{tabularx}
\end{table}

\begin{tcolorbox}[gotchabox]
\begin{itemize}
    \item \textbf{Swara-Bhakti in Action:} Both \deva{परताप} (Parataap) and \deva{परसिद्ध} (Parasiddha) exhibit pure Awadhi Swara-Bhakti (anaptyxis). The strict Sanskrit conjuncts \deva{प्रताप} (Prataapa) and \deva{प्रसिद्ध} (Prasiddha) are broken open by inserting a short `a` vowel between $P$ and $R$. This makes the words musically softer and easier to sing while perfectly occupying the required musical beats.
\end{itemize}
\end{tcolorbox}

\newpage
